% !TeX program = xelatex
\documentclass[11pt,a4paper]{article}
\usepackage[margin=23mm]{geometry}
\usepackage{fontspec}
\setmainfont{Latin Modern Roman}
\usepackage{amsmath,amssymb,mathtools}
\usepackage{xcolor,enumitem,booktabs,longtable,array,xurl,hyperref}
\definecolor{muted}{HTML}{53636C}
\hypersetup{colorlinks=true,urlcolor=blue,linkcolor=blue}
\setlength{\parindent}{0pt}
\setlength{\parskip}{5pt}
\newcommand{\R}{\mathbb R}
\newcommand{\E}{\mathbb E}
\newcommand{\N}{\mathbb N}
\newcommand{\ind}{\mathbf 1}
\newcommand{\Var}{\operatorname{Var}}
\newcommand{\Cov}{\operatorname{Cov}}
\newcommand{\supp}{\operatorname{supp}}
\DeclareMathOperator*{\argmin}{arg\,min}
\DeclareMathOperator*{\argmax}{arg\,max}
\newcommand{\entrymeta}[3]{{\sffamily\small\color{muted}\textbf{\texttt{#1}}\quad|\quad #2\par #3\par}}
\newcommand{\Problemref}[1]{\texttt{#1}}
\newcommand{\JELref}[2]{\texttt{#2} (JEL \texttt{#1})}
\begin{document}
\section*{Sequential adaptive experiments with uncertain causal structure}
\textbf{Problem ID:} \texttt{C90-42}\quad \textbf{JEL:} \texttt{C90}\par
% BEGIN ECONOMICS STATEMENT
\label{problem:OP-0480}
\label{beyond:CI-01}
\entrymeta{CI-01}{Econometric theory and causal inference}{Research agenda formalization}
\subsection*{Definitions and assumptions}
Fix a finite horizon $K\geq1$, finite action spaces $\mathcal W_k$, and an observed longitudinal record
\[
 O=(X,W_1,L_1,\ldots,W_{K-1},L_{K-1},W_K,Y),\qquad 0\leq Y\leq1.
\]
For $K=1$ there are no intermediate $L$ variables. The measurable spaces of $X$ and $L_k$ are specified inputs; $L_k$ is observed after $W_k$ and before $W_{k+1}$. Write $H_k=(X,W_1,L_1,\ldots,W_{k-1},L_{k-1})$, with $H_1=X$. Let $\mathcal G$ be a specified finite collection of causal directed acyclic graphs, possibly containing latent variables, sharing these observed node labels and this observation order. The specification of $\mathcal M(\mathcal G)$ includes the latent variable spaces, graph-to-structural-model interpretation, temporal compatibility, structural equations and permitted exogenous dependence restrictions. No independent-error or unrestricted-error convention is inferred merely from the graph notation.

Index structural models by $m\in\mathcal M(\mathcal G)$. A target regime $g=(g_1,\ldots,g_K)$ must specify measurable maps $g_k$ from the space of $H_k$ to $\mathcal W_k$, together with their intervention semantics. Constant maps $g_k\equiv w_k$ include the stated static targets. Each model must give unique measurable counterfactuals $Y^g$ and satisfy $\E_m|Y^g|<\infty$ for every target regime under consideration; write $\theta_g(m)=\E_m[Y^g]$. Any additional counterfactual range or uniform moment restriction is part of the declared structural class, not a consequence asserted here of the observed bound on $Y$.

Across $n$ episodes with independent structural errors before assignment, the design may choose assignment kernels $q_{i,k}$ using previous episodes and the current observed history. Observed episodes can therefore be dependent. Assignments use external randomness independent of structural errors conditional on that history. Specify which treatment coordinates can be manipulated and their allowable probabilities; latent variables remain unrestricted except by the declared structural class.

For each $n\geq1$, the nonempty allowable family $\mathcal Q_n$ is an input of the agenda. Its specification gives the information available before each assignment, including the available feedback from earlier episodes; the kernels governing manipulated coordinates; the structural equations retained for natural coordinates; and the ordering and joint assignment semantics when only part of an action is manipulated. It also gives all support, probability, budget or other restrictions and their possible dependence on $n$. A design is a known rule chosen without knowledge of $m$, although it may adapt to its permitted information. The recorded transcript $D_n$ and the estimator's information, including access to design rules, assignment probabilities and any recorded design randomization, are specified with $\mathcal Q_n$. The episode sampling specification must additionally state the population defining $\theta_g(m)$ and its relation to the experimental episodes; independence of structural errors alone does not choose identical episode laws, invariant mechanisms or absence of cross-episode causal effects.
\subsection*{Formal question}
For each fully specified structural and intervention family, let $P_m^{q,n}$ be the law of $D_n$ under model $m$ and design $q\in\mathcal Q_n$. Characterize when
\[
 \bigl[P_m^{q,n}=P_{m'}^{q,n}\text{ for every }n\geq1
       \text{ and every }q\in\mathcal Q_n\bigr]
 \quad\Longrightarrow\quad \theta_g(m)=\theta_g(m')
 \qquad(m,m'\in\mathcal M(\mathcal G)).
\]
Here all observable adaptive experiments means all these finite experiments. When identification holds, determine the minimax mean squared risk
\[
 R_n=\inf_{q\in\mathcal Q_n,\,\widehat\theta_n}
       \sup_{m\in\mathcal M(\mathcal G)}
       \E_{P_m^{q,n}}[(\widehat\theta_n(D_n)-\theta_g(m))^2],
\]
where estimators are measurable with respect to their declared information and the design-estimator pair is selected before the unknown model. Identification alone does not assert a vanishing uniform risk. Rate statements specify which class parameters are fixed, which vary with $n$, and the permitted dependence of constants on horizon, graph uncertainty, action-space sizes and intervention restrictions.

When identification fails, determine the sharp identified set
\[
 \Theta_g(Q)=\{\theta_g(m):m\in\mathcal M(\mathcal G),\
       P_m^{q,n}=Q_{n,q}\text{ for every }n\geq1,\ q\in\mathcal Q_n\},
\]
where $Q=(Q_{n,q})_{n\geq1,\,q\in\mathcal Q_n}$ is an indexed family of observable laws. The question concerns attainable families; an incompatible family has an empty identified set. For an implemented $q\in\mathcal Q_n$ and an attainable transcript law $A$, its corresponding identified set is separately
\[
 \Theta_g^{q,n}(A)=\{\theta_g(m):m\in\mathcal M(\mathcal G),\ P_m^{q,n}=A\}.
\]
For any implemented design, seek informative confidence sets for this corresponding identified set. The coverage event, uniformity, fixed-time or time-uniform requirement, and accuracy criterion relative to the identified set must be specified; the scalar coverage convention below does not choose them. Positive identified-set diameter is not itself inferential error. Establish matching lower bounds in terms of intervention support, graph uncertainty, and horizon for the declared classes and statistical criteria.
\subsection*{Scope and status}
Complete randomization of every needed action with adequate support already resolves many identification questions. The research target is the specified restricted-intervention class and its sharp statistical limits; ignorance of confounders cannot by itself justify point identification. This is a source-dated research agenda refinement.
\subsection*{References for the source of the problem}
Carlos Cinelli, Avi Feller, Guido Imbens, Edward Kennedy, Sara Magliacane, and Jose Zubizarreta, \emph{Challenges in Statistics: A Dozen Challenges in Causality and Causal Inference}, arXiv:2508.17099v1, 23 August 2025, Section 3.1.3, paragraph \emph{Sequential treatments and increased adaptivity}. \url{https://arxiv.org/html/2508.17099v1}.

\subsection*{Common conventions: Source mathematical conventions}

\textbf{Well-defined objects.}
All probability spaces and measurable variables are specified or inherited
from a local statistical model. Expectations used as finite targets require
the stated integrability. A decision rule or estimator is measurable with
respect to its available information; randomized rules may use an auxiliary
independent random variable. Norms without a subscript are Euclidean in
finite-dimensional spaces. An optimizer is requested only when attainment
is assured; otherwise the question concerns the optimal value and
$\varepsilon$-optimal admissible rules for every $\varepsilon>0$.

\textbf{Identification.}
A structural model $m\in\mathcal M$ implies an observable law
$P_m^{\rm obs}$ and target $\theta(m)$. The target is point identified on
$\mathcal M$ if
\[
 P_{m_1}^{\rm obs}=P_{m_2}^{\rm obs}
 \quad\Longrightarrow\quad\theta(m_1)=\theta(m_2)
 \qquad(m_1,m_2\in\mathcal M).
\]
When this fails, the sharp identified set is
\[
 \Theta(P;\mathcal M)=\{\theta(m):m\in\mathcal M,
                                  P_m^{\rm obs}=P\}.
\]
Specifying an intervention or estimand does not establish identification.
Positivity, exclusion, causal sufficiency, measurement restrictions, and
transport assumptions are substantive local inputs, never automatic
consequences of notation.

\textbf{Minimax risk and nontrivial inference.}
For a statistical experiment with data law $P^{(n)}$, target $\theta(P)$,
and nonnegative loss $\ell$, define
\[
 \mathcal R_n(\mathcal P,\ell)
   =\inf_{\widehat\theta_n}\sup_{P\in\mathcal P}
      \E_{P^{(n)}}\ell(\widehat\theta_n,\theta(P)).
\]
The infimum is over measurable estimators. A sharp rate is a sequence
$r_n$ for which $0<c\leq C<\infty$, independent of $n$, satisfies
$c r_n\leq\mathcal R_n\leq C r_n$ for all sufficiently large $n$.
Squared-error parametric risk is of order $n^{-1}$; the corresponding
root-mean-squared error is of order $n^{-1/2}$. These different conventions
are distinguished locally. A uniform asymptotic confidence region $C_n$
must satisfy
\[
 \liminf_{n\to\infty}\inf_{P\in\mathcal P}
     \Pr_{P^{(n)}}\{\theta(P)\in C_n\}\geq 1-\alpha,
 \qquad 0<\alpha<1,
\]
together with an informative diameter or expected-width requirement.
Coverage alone permits the vacuous whole parameter space. Testing questions
similarly impose both size control and power against a declared separated
alternative class.

\textbf{Binary causal baseline.}
When an entry explicitly invokes this baseline, one observes independent
copies of $O=(X,W,Y)$, with $W\in\{0,1\}$ and potential outcomes $Y(0),Y(1)$.
Assume consistency $Y=Y(W)$, no interference, conditional exchangeability
$(Y(0),Y(1))\perp W\mid X$, and overlap
$\epsilon\leq e(X)\leq1-\epsilon$ almost surely for a fixed
$0<\epsilon<1/2$, where
\[
 e(x)=\Pr(W=1\mid X=x),\qquad
 m_w(x)=\E[Y\mid W=w,X=x],\qquad
 \tau(x)=m_1(x)-m_0(x).
\]
These assumptions identify conditional mean effects almost everywhere
under the covariate law. Pointwise or uniform effects require appropriate
regularity and a specified support. Continuous treatment, longitudinal
data, adaptive experiments, and network interference require their own
assumptions in place of this baseline. Bounded outcomes, densities, and
smoothness are additional assumptions only where stated.

\textbf{Function classes.}
On $[0,1]^d$, $\mathcal H_d^s(L)$ is the isotropic Hölder ball of radius
$L>0$: put $k=\lceil s\rceil-1$ and $\gamma=s-k\in(0,1]$, bound derivatives
through order $k$, and require the order-$k$ derivatives to be
$\gamma$-Hölder with constant at most $L$. Thus integer orders use a
Lipschitz final derivative convention. The class
$\operatorname{BV}_d(L)$ consists of functions $f\in L^1((0,1)^d)$ whose
distributional derivative is a finite vector measure and for which
$\|f\|_{L^1}+|Df|((0,1)^d)\leq L$.
An $L^\infty$ bound or a particular pointwise version is imposed separately
when needed. Sparse and neural-network classes require a declared
dictionary or architecture and coefficient bounds.

An anisotropic H\"older class with declared block smoothness indices means
coordinatewise H\"older regularity in each block, uniformly in the other
blocks, with all corresponding derivative and H\"older bounds at most
the stated radius. Derivative orders and the integer-order convention in
each block are those just given. Mixed-derivative bounds are additional
restrictions only when explicitly stated.

\textbf{Decision and computational questions.}
Policy regret compares admissible feasible policies under the same law and
constraints. In computational entries, finite data and thresholds have
rational encodings; running time is measured in their total bit length.
Real-valued equilibrium search uses the explicitly declared algebraic or
FIXP encoding. An approximation ratio for maximization is written
$\operatorname{OPT}/\operatorname{ALG}$ and for minimization as
$\operatorname{ALG}/\operatorname{OPT}$, with zero optima and infeasible
instances handled locally. Historical approximation bounds are not
automatically current bounds.
% END ECONOMICS STATEMENT
\end{document}
